%===============================================================================
% STANDALONE EXAMPLE: Using Template from Project Root
%===============================================================================
% This example shows how to use the template when it's cloned as a subfolder
% in your project root (the recommended approach for most projects).
%
% Directory structure:
%   your_project/
%   ├── gimmatek_report_latex_template/  ← cloned template
%   │   ├── src/
%   │   │   ├── gimmatek-report-template.cls
%   │   │   └── gimmatek-report-frontmatter.sty
%   │   └── assets/
%   ├── standalone_example.tex  ← this file (copy to your project root)
%   └── chapters/
%       ├── chapter1.tex
%       └── chapter2.tex
%===============================================================================

\documentclass[12pt,a4paper]{gimmatek-report-template}

% Point to template packages in subfolder
\makeatletter
\def\input@path{{gimmatek_report_latex_template/src/}}
\makeatother
\RequirePackage{gimmatek-report-frontmatter}

%-------------------------------------------------------------------------------
% Document Metadata
%-------------------------------------------------------------------------------
\doctitle{Standalone Project Example}
\doctitlepage{Standalone Project\\Example Document}
\docsubtitle{Using Template from Project Root}
\docversion{V1.0}
\doctype{Technical Documentation}

%-------------------------------------------------------------------------------
% Abstract
%-------------------------------------------------------------------------------
\docabstract{%
This example demonstrates using the Gimmatek LaTeX Template from your project
root directory. This is the recommended setup for most projects as it keeps
the template self-contained and version-controlled within your project.

Key benefits of this approach:
\begin{itemize}
\item Template is version-locked to your project
\item No system-wide installation required
\item Works offline and in CI/CD environments
\item Easy to share with collaborators
\end{itemize}
}

%===============================================================================
% DOCUMENT BODY
%===============================================================================
\begin{document}

%-------------------------------------------------------------------------------
% Front Matter
%-------------------------------------------------------------------------------
\frontmatter
\makefrontmatter

%-------------------------------------------------------------------------------
% Main Content
%-------------------------------------------------------------------------------
\mainmatter

\chapter{Introduction}

This document shows the recommended way to use the Gimmatek LaTeX Template:
clone it into your project as a subfolder.

\section{Setup Instructions}

To use this approach in your own project:

\begin{enumerate}
\item Clone the template into your project:
\begin{verbatim}
cd your_project/
git clone https://github.com/gimmatekmac05/\
  gimmatek_report_latex_template.git
\end{verbatim}

\item Copy this file to your project root

\item Update the metadata section with your document details

\item Create a \texttt{chapters/} folder for your content

\item Compile with XeLaTeX:
\begin{verbatim}
xelatex standalone_example.tex
xelatex standalone_example.tex
\end{verbatim}
\end{enumerate}

\section{Path Configuration}

The key difference from the examples directory version is the path configuration:

\begin{verbatim}
\makeatletter
\def\input@path{{gimmatek_report_latex_template/src/}}
\makeatother
\end{verbatim}

This tells LaTeX to look for packages in the \texttt{gimmatek\_report\_latex\_template/src/}
subdirectory relative to your document.

\chapter{Features}

\section{Automatic Front Matter}

The template automatically generates all front matter elements:

\begin{table}[htbp]
\centering
\begin{tabularx}{\textwidth}{l X}
\toprule
\textbf{Element} & \textbf{Description} \\
\midrule
Abstract & Chinese-titled abstract page (摘要) \\
Table of Contents & Automatic TOC with configurable depth \\
List of Tables & Generated if document contains tables \\
List of Figures & Generated if document contains figures \\
Page Numbering & Roman numerals (front matter) → Arabic (main) \\
\bottomrule
\end{tabularx}
\caption{Automatically Generated Elements}
\label{tab:elements}
\end{table}

\section{Customization Options}

You can customize the front matter by setting flags before calling
\texttt{\textbackslash makefrontmatter}:

\begin{verbatim}
% Disable specific elements
\includetocfalse     % Skip table of contents
\includelotfalse     % Skip list of tables
\includeloffalse     % Skip list of figures
\includeabstractfalse % Skip abstract
\end{verbatim}

\section{Logo Configuration}

The template includes a configurable logo system with automatic fallback:

\begin{verbatim}
% Use custom logo (optional)
\logopath{path/to/your/logo.png}
\end{verbatim}

If the logo file doesn't exist, the template will display a placeholder
instead of failing.

\chapter{Multi-File Projects}

\section{Using Separate Chapter Files}

For larger projects, split content into separate files:

\begin{verbatim}
% In your main .tex file:
\mainmatter
\input{chapters/chapter1.tex}
\input{chapters/chapter2.tex}
\input{chapters/chapter3.tex}
\end{verbatim}

\section{Example Chapter Structure}

Create \texttt{chapters/chapter1.tex}:

\begin{verbatim}
\chapter{First Chapter}

\section{Introduction}

Content goes here...
\end{verbatim}

No \texttt{\textbackslash begin\{document\}} or
\texttt{\textbackslash end\{document\}} in chapter files!

\chapter{Compilation}

\section{Using Makefile}

For convenience, create a \texttt{Makefile}:

\begin{verbatim}
TEXINPUTS := .:./gimmatek_report_latex_template/src//:
export TEXINPUTS

all: standalone_example.pdf

standalone_example.pdf: standalone_example.tex
    xelatex standalone_example.tex
    xelatex standalone_example.tex

clean:
    rm -f *.aux *.log *.toc *.lot *.lof *.out

.PHONY: all clean
\end{verbatim}

Then simply run: \texttt{make}

\section{Environment Variable Method}

Alternatively, set \texttt{TEXINPUTS} in your shell:

\begin{verbatim}
export TEXINPUTS=.:./gimmatek_report_latex_template/src//:
xelatex standalone_example.tex
\end{verbatim}

\chapter{Version Control}

\section{Git Integration}

You have two options for version control:

\subsection{Option 1: Git Submodule (Recommended)}

\begin{verbatim}
git submodule add \
  https://github.com/gimmatekmac05/\
  gimmatek_report_latex_template.git
\end{verbatim}

Benefits:
\begin{itemize}
\item Template updates tracked separately
\item Easy to update template version
\item Cleaner project history
\end{itemize}

\subsection{Option 2: Direct Clone}

\begin{verbatim}
git clone https://github.com/gimmatekmac05/\
  gimmatek_report_latex_template.git
\end{verbatim}

Benefits:
\begin{itemize}
\item Simpler setup
\item Self-contained project
\item No submodule complexity
\end{itemize}

\chapter{Conclusion}

This standalone approach provides the best balance of:
\begin{itemize}
\item Ease of use
\item Portability
\item Version control
\item Team collaboration
\end{itemize}

Start your own project by copying this file and customizing it for your needs!

\end{document}

%===============================================================================
% QUICK START
%===============================================================================
%
% 1. Clone template into your project:
%    git clone https://github.com/gimmatekmac05/\
%      gimmatek_report_latex_template.git
%
% 2. Copy this file to project root
%
% 3. Compile:
%    xelatex standalone_example.tex
%    xelatex standalone_example.tex
%
% 4. Done! Edit and rebuild as needed.
%
%===============================================================================
